\documentclass[10pt,a4paper]{amsart}
\usepackage[margin=19mm]{geometry}
\usepackage{amsmath,amssymb,mathtools}
\usepackage[scr=rsfs,cal=euler]{mathalfa}
\usepackage{xfrac}
\usepackage[unicode,hidelinks,bookmarks=false]{hyperref}
\newcommand{\ie}{i.e.\ }
\newcommand{\cf}{cf.\ }
\newcommand{\resp}{respectively\ }
\newcommand{\Subsection}{\subsection}
\providecommand{\nobreakdash}{\nobreak\hbox{-}}
\setlength{\emergencystretch}{2em}
\multlinegap0pt
\usepackage{amssymb}
\usepackage{mathtools}
\usepackage{xcolor}
\usepackage{algorithm}\usepackage{algpseudocode}
\newtheorem{theorem}{Theorem}
\newcommand{\E}[2][]{\mathbb{E}_{#1}\!\left[\,#2\,\right]}
\newcommand{\ind}{\stackrel{\text{ind}}{\sim}}
\newcommand{\msehat}{\widehat{\mathrm{MSE}}_{\epsilon}}
\newcommand{\msethin}{\mathrm{MSE}_{\epsilon}}
\newcommand{\thetatrain}{\hat{\theta}^{(1)}}           % Training estimator
\newcommand{\ytest}{y^{(2)}}                            % Test data
\newcommand{\ytrain}{y^{(1)}}                           % Training data
\title{On Data Thinning for Model Validation in Small Area Estimation}
\author{Sho Kawano, Paul A. Parker, Zehang Richard Li}
\date{Source: arXiv:2604.04141v3 (CC BY 4.0)}
\begin{document}
\maketitle

\noindent\emph{Source and licence.} On Data Thinning for Model Validation in Small Area Estimation. Version arXiv:2604.04141v3 (CC BY 4.0), distributed by arXiv under the Creative Commons Attribution 4.0 International licence, http://creativecommons.org/licenses/by/4.0/, as recorded in the arXiv licence metadata for this exact version and shown on the abstract page https://arxiv.org/abs/2604.04141v3. Full licence notice: \texttt{licenses/arxiv-2604.04141-CC-BY-4.0.txt}.\par\smallskip

\noindent\textit{Editorial scope.} Counted unit: the article's theorem thm:unbiased\_mse (source0.tex lines 333--340). The article's complete original proof of it is delivered with it (source0.tex lines 342--379, one \texttt{proof} environment of 38 lines). No mathematics is written, repaired or re-derived: the statement and the proof are the article's own lines, in the article's own notation, and no retained line is substituted or re-flowed.  No further source-local context is needed: every cross-reference of every retained line resolves inside the two delivered spans.  Nothing was omitted from any retained span, so the package carries no omission record.  No retained line contains a citation, so the wrapper carries no bibliography and no bibliography entry.  No retained line contains a figure environment, an image-inclusion command or drawing code, so no image is delivered and the package declares no asset dependency.  Editorial formatting only, in addition: an AMS \texttt{amsart} preamble that restates the article's own declarations and macros used by the retained lines, namely the environments \texttt{theorem} on separate counters, together with the standard packages those lines need.  The retained environments print in this document as Theorem 1 in order of appearance, whereas the article prints the same items with the numbering of its own full document.  The complete native source of the article as distributed in the arXiv e-print for this exact version is bundled unchanged as \texttt{sources/197863/source0.tex}.
\begin{theorem}[Unbiased MSE estimation]
\label{thm:unbiased_mse}
The estimator $\msehat$ is unbiased for the thinned-data oracle MSE:
\[
\E{\msehat} = \msethin,
\]
where the expectation is taken over the joint distribution of $(y^{(1)}, y^{(2)})$, unconditional on $y$.
\end{theorem}

\begin{proof}
We take the expectation over the joint distribution of $\ytrain$ and $\ytest$ to get
\begin{align*}
    \E{\msehat}
& = \frac{1}{m}\sum_{i=1}^m
   \E[\ytrain, \ytest]{ \left( \thetatrain_i - \tfrac{1}{1-\epsilon}\ytest_i \right)^{\!2}
         - \tfrac{d_i}{1-\epsilon}} \\
&=  \frac{1}{m}\sum_{i=1}^m\E[\ytrain]{\E[\ytest]{\left( \thetatrain_i - \tfrac{1}{1-\epsilon}\ytest_i \right)^{\!2}} - \tfrac{d_i}{1-\epsilon}},
\end{align*}
where the second equality results from the marginal independence of $\ytrain$ and $\ytest$, unconditional on $y$.

Fix an area $i$ and define
\[
\delta_i := \thetatrain_i - \theta_i,
\qquad
\eta_i := \tfrac{1}{1-\epsilon}\ytest_i - \theta_i .
\]
Note that $\delta_i$ is only random with respect to $\ytrain$, while $\eta_i \ind N(0, d_i/(1-\epsilon))$ is random with respect to $\ytest$.

Thus we have
\[
\E[\ytest]{(\delta_i-\eta_i)^2}
=
\delta_i^2 - 2 \delta_i \cdot \E[\ytest]{\eta_i} +  \E[\ytest]{\eta_i^2}
=
\delta_i^2 + \tfrac{d_i}{1-\epsilon}.
\]
Substituting into the definition of $\msehat$ and taking $\E[\ytrain]{\cdot}$ gives
\begin{align*}
    \E{\msehat}
&=  \frac{1}{m}\sum_{i=1}^m\E[\ytrain]{\E[\ytest]{(\delta_i-\eta_i)^2} - \tfrac{d_i}{1-\epsilon}}\\
&=  \frac{1}{m}\sum_{i=1}^m \E[\ytrain]{\delta_i^2 + \tfrac{d_i}{1-\epsilon} - \tfrac{d_i}{1-\epsilon}}  \\
&=
\frac{1}{m}\sum_{i=1}^m \E[\ytrain]{\left( \thetatrain_i - \theta_i \right)^2}
=
\msethin.
\end{align*}
\end{proof}

\end{document}
